Полученные выводы следует интерпретировать с учётом ограничений
бенчмарка. Для каждой комбинации задачи, метода и бюджета использовано
только 10 целевых наблюдений. Ограничение связано с вычислительной
стоимостью экспериментов; увеличение числа posterior-сэмплов для каждого
наблюдения не заменяет увеличение числа независимых целевых наблюдений.
Поэтому агрегированные показатели калибровки, включая TARP и SBC, имеют
ограниченную статистическую мощность и носят ориентировочный характер.

Для SIR и StochVol нет аналитического эталонного posterior, поэтому
качество оценивалось по калибровке и posterior-predictive проверкам.
Хорошая калибровка сама по себе не означает, что оценённая плотность
совпадает с истинным posterior. Кроме того, в StochVol в бенчмарке
используется пятикомпонентная сводка доходностей вместо полного временного
ряда; выводы относятся к posterior параметров при условии этой сводки.

Отдельное ограничение связано с тем, что нейросетевые методы могут
генерировать параметры вне допустимой области модели. Для SIR проверяется
условие $0.05\leq\gamma\leq0.5$ и $0.5\leq\beta/\gamma\leq5$, а для
StochVol --- $-1<\phi<1$ и $\sigma>0$. При одном раунде доля сэмплов
SFMPE, попавших в допустимую область, для бюджетов 100, 1000 и 10000
составляет
$0.523$, $0.727$ и $0.876$ на SIR и $0.392$, $0.490$ и $0.905$ на
StochVol. Для SMC-ABC соответствующие значения равны $0.739$, $0.842$ и
$0.869$ на SIR, а также $0.610$, $0.757$ и $0.781$ на StochVol. В
сохранённых результатах SNPE-C все сэмплы проходят ту же проверку
допустимой области. Поэтому показатели калибровки SFMPE и SMC-ABC следует
рассматривать вместе с этими долями: значение ATC само по себе не
отражает, сколько сэмплов нарушают ограничения модели.

Наконец, измеренное время выполнения зависит от аппаратного устройства
и фоновой нагрузки, а запуски бенчмарка выполнялись с автоматическим
выбором устройства. Поэтому приведённые значения времени дают лишь
ориентировочное сравнение. Бюджет симуляций учитывает симуляции,
использованные для инференса, но не posterior-сэмплирование и
дополнительные симуляции для диагностик.